\color{unidarkgreen}{\section[(обязательное) Уставки функций РЗА]{(обязательное)\\Уставки функций РЗА}\label{app:set}}
\color{black}



%%%%%%%%%%%%%%%%%%%%%%%%%%%%%%%%%%%%%%%%%%%%%%%%%%%%%%%%%% База %%%%%%%%%%%%%%%%%%%%%%%%%%%%%%%%%%%%%%%%%%%%%%%%%%%%%%%%%%%%%%%%%%
\par
В таблице \ref{obsh:tbl1} приведены номинальные параметры сети и основные настройки АРН.

\footnotesize
\begin{longtable}{|>{\centering\arraybackslash}m{5.3cm}|>{\centering\arraybackslash}m{3.3cm}|>{\centering\arraybackslash}m{4.2cm}|>{\centering\arraybackslash}m{1.8cm}|>{\centering\arraybackslash}m{1cm}|}
\caption{Основные настройки и номинальные параметры сети\hfill\vspace{-0.5\baselineskip}}\label{obsh:tbl1}\\ 
\hline
\rowcolor{gray!30}
Параметр (Параметр на ИЧМ) & Условное обозначение на схеме & Значение/ Диапазон & Единица измерения & Шаг \\ 
\hline
\endfirsthead
\caption*{\hspace{3pt}\emph{Продолжение таблицы \ref{obsh:tbl1}\hfill\vspace{-0.5\baselineskip}}} \\ % сделано по ГОСТ 2.105 п.6.8.7
\hline
\rowcolor{gray!30}
Параметр (Параметр на ИЧМ) & Условное обозначение на схеме & Значение/ Диапазон & Единица измерения & Шаг \\ 
%\hline
\endhead
%\hline
\endfoot
\hline
\endlastfoot
%==+t1*LV1TCTR1|ARNT_ttccMEAS
\centering Схема подключения вторичных обмоток ТН 1 секции (Подкл. ЦН ТН1)  & \centering Подкл. ЦН ТН1 & \centering Ua / Ub / Uc \\ Uab / Ubc / 3U0 & \centering -- & \centering \arraybackslash -- \\
\hline
\centering Схема подключения вторичных обмоток ТН 2 секции (Подкл. ЦН ТН2)  & \centering Подкл. ЦН ТН2 & \centering Ua / Ub / Uc \\ Uab / Ubc / 3U0 & \centering -- & \centering \arraybackslash -- \\
\hline
\centering Измеряемая фаза тока вводного выключателя 1 секции (Фаза тока Iвв1)  & \centering Фаза тока Iвв1 & \centering А \\ В \\ С & \centering -- & \centering \arraybackslash -- \\
\hline
\centering Измеряемая фаза тока секционного выключателя 1 секции (Фаза тока Iсв1)  & \centering Фаза тока Iсв1 & \centering А \\ В \\ С & \centering -- & \centering \arraybackslash -- \\
\hline
\centering Измеряемая фаза тока вводного выключателя 2 секции (Фаза тока Iвв2)  & \centering Фаза тока Iвв2 & \centering А \\ В \\ С & \centering -- & \centering \arraybackslash -- \\
\hline
\centering Измеряемая фаза тока секционного выключателя 2 секции (Фаза тока Iсв2)  & \centering Фаза тока Iсв2 & \centering А \\ В \\ С & \centering -- & \centering \arraybackslash -- \\
\hline
\centering Номинальное напряжение 1 секции (Uном 1с)  & \centering Uном 1с & \centering 100,00...999999,99 & \centering В & \centering \arraybackslash 0,01 \\
\hline
\centering Номинальное напряжение 2 секции (Uном 2с)  & \centering Uном 2с & \centering 100,00...999999,99 & \centering В & \centering \arraybackslash 0,01 \\
\hline
\centering Номинальный ток ввода 1 секции (Iном вв1)  & \centering Iном вв1 & \centering 1,00...99999,99 & \centering А & \centering \arraybackslash 0,01 \\
\hline
\centering Номинальный ток ввода 2 секции (Iном вв2)  & \centering Iном вв2 & \centering 1,00...99999,99 & \centering А & \centering \arraybackslash 0,01 \\
\hline
\centering Номинальный ток обмотки трансформатора с РПН (Iном РПН)  & \centering Iном РПН & \centering 1,00...99999,99 & \centering А & \centering \arraybackslash 0,01 \\
\hline
\centering Активное сопротивление сети 1 секции для линейной компенсации (Rлинии 1секции)  & \centering Rлинии 1секции & \centering 0,00...99999,99 & \centering Ом & \centering \arraybackslash 0,01 \\
\hline
\centering Реактивное сопротивление сети 1 секции для линейной компенсации (Xлинии 1секции)  & \centering Xлинии 1секции & \centering 0,00...99999,99 & \centering Ом & \centering \arraybackslash 0,01 \\
\hline
\centering Активное сопротивление сети 2 секции для линейной компенсации (Rлинии 2секции)  & \centering Rлинии 2секции & \centering 0,00...99999,99 & \centering Ом & \centering \arraybackslash 0,01 \\
\hline
\centering Реактивное сопротивление сети 2 секции для линейной компенсации (Xлинии 2секции)  & \centering Xлинии 2секции & \centering 0,00...99999,99 & \centering Ом & \centering \arraybackslash 0,01 \\
\hline
%===t1
\hline
\end{longtable}
\normalsize

%%%%%%%%%%%%%%%%%%%%%%%%%%%%%%%%%%%%%%%%%%%%%%%%%%%%%%%%%% АРН %%%%%%%%%%%%%%%%%%%%%%%%%%%%%%%%%%%%%%%%%%%%%%%%%%%%%%%%%%%%%%%%%%
\par
В таблице \ref{kntr:tbl1} приведены параметры, необходимые для настройки автоматики регулирования напряжения.

\footnotesize
\begin{longtable}{|>{\centering\arraybackslash}m{5.3cm}|>{\centering\arraybackslash}m{3.3cm}|>{\centering\arraybackslash}m{4.2cm}|>{\centering\arraybackslash}m{1.8cm}|>{\centering\arraybackslash}m{1cm}|}
\caption{Параметры для настройки АРН\hfill\vspace{-0.5\baselineskip}}\label{kntr:tbl1}\\ 
\hline
\rowcolor{gray!30}
Параметр (Параметр на ИЧМ) & Условное обозначение на схеме & Значение/ Диапазон & Единица измерения & Шаг \\ 
\hline
\endfirsthead
\caption*{\hspace{3pt}\emph{Продолжение таблицы \ref{kntr:tbl1}\hfill\vspace{-0.5\baselineskip}}} \\ % сделано по ГОСТ 2.105 п.6.8.7
\hline
\rowcolor{gray!30}
Параметр (Параметр на ИЧМ) & Условное обозначение на схеме & Значение/ Диапазон & Единица измерения & Шаг \\ 
%\hline
\endhead
%\hline
\endfoot
\hline
\endlastfoot
%==+t1*ATCC1|ARNT_TTCC
\centering Источник сигнала положения РПН (Датчик\_РПН)  & \centering Датчик\_РПН & \centering Внутренний счетчик \\ Токовая петля 4(0)-20мА \\ GOOSE \\ BCD-код & \centering -- & \centering \arraybackslash -- \\
\hline
\centering Первоначальное значение счетчика положений РПН (Нач\_полож\_РПН)  & \centering Нач\_полож\_РПН & \centering 1...40 & \centering -- & \centering \arraybackslash 1 \\
\hline
\centering Номер крайнего нижнего положения РПН (Мин\_положение\_РПН)  & \centering Мин\_положение\_РПН & \centering 1...40 & \centering -- & \centering \arraybackslash 1 \\
\hline
\centering Номер крайнего верхнего положения РПН (Макс\_положение\_РПН)  & \centering Макс\_положение\_РПН & \centering 1...40 & \centering -- & \centering \arraybackslash 1 \\
\hline
\centering Направление счета положений (Направление\_счета)  & \centering Направление\_счета & \centering Прямо \\ Обратно & \centering -- & \centering \arraybackslash -- \\
\hline
\centering Сигнал токовой петли, соответствующий крайнему нижнему положению РПН (Iмин\_ТП\_РПН)  & \centering Iмин\_ТП\_РПН & \centering 0...4000 & \centering мкА & \centering \arraybackslash 1 \\
\hline
\centering Сигнал токовой петли, соответствующий крайнему верхнему положению РПН (Iмакс\_ТП\_ РПН)  & \centering Iмакс\_ТП\_ РПН & \centering 5000...20000 & \centering мкА & \centering \arraybackslash 1 \\
\hline
\centering Положение привода, далее которого блокируются автоматические команды прибавить (Верхн\_разреш\_полож)  & \centering Верхн\_разреш\_полож & \centering 1...40 & \centering -- & \centering \arraybackslash 1 \\
\hline
\centering Положение привода, далее которого блокируются автоматические команды убавить (Нижн\_разреш\_полож)  & \centering Нижн\_разреш\_полож & \centering 1...40 & \centering -- & \centering \arraybackslash 1 \\
\hline
\centering Промежуточное 1 положение РПН (Промежут\_1\_полож)  & \centering Промежут\_1\_полож & \centering 0...40 & \centering -- & \centering \arraybackslash 1 \\
\hline
\centering Промежуточное 2 положение РПН (Промежут\_2\_полож)  & \centering Промежут\_2\_полож & \centering 0...40 & \centering -- & \centering \arraybackslash 1 \\
\hline
\centering Промежуточное 3 положение РПН (Промежут\_3\_полож)  & \centering Промежут\_3\_полож & \centering 0...40 & \centering -- & \centering \arraybackslash 1 \\
\hline
\centering Промежуточное 4 положение РПН (Промежут\_4\_полож)  & \centering Промежут\_4\_полож & \centering 0...40 & \centering -- & \centering \arraybackslash 1 \\
\hline
\centering Промежуточное 5 положение РПН (Промежут\_5\_полож)  & \centering Промежут\_5\_полож & \centering 0...40 & \centering -- & \centering \arraybackslash 1 \\
\hline
\centering Центр зоны нечувствительности (Uпод)  & \centering Uпод & \centering 0,85...1,45 & \centering о.е. & \centering \arraybackslash 0,01 \\
\hline
\centering Ширина зоны нечувствительности (dUпод)  & \centering dUпод & \centering 0,01...0,20 & \centering о.е. & \centering \arraybackslash 0,01 \\
\hline
\centering Уровень напряжения, при превышении которого автоматические команды на снижение формируются с ускоренной выдержкой времени T3 (U\verb|>>|сраб)  & \centering U\verb|>>|сраб & \centering 1,00...2,00 & \centering о.е. & \centering \arraybackslash 0,01 \\
\hline
\centering Величина тока перегрузки РПН (Iрпн>блк)  & \centering Iрпн>блк & \centering 1,00...2,00 & \centering о.е. & \centering \arraybackslash 0,01 \\
\hline
\centering Блокировка от контролируемой секции (БлкКонтрСекц)  & \centering БлкКонтрСекц & \centering Не предусмотрено \\ Предусмотрено & \centering -- & \centering \arraybackslash -- \\
\hline
\centering Величина тока перегрузки для запрета авторегулирования на повышение (I>блк)  & \centering I>блк & \centering 1,000...2,000 & \centering о.е. & \centering \arraybackslash 0,001 \\
\hline
\centering Величина напряжения НП для запрета авторегулирования на повышение (3U0блк)  & \centering 3U0блк & \centering 0,00...2,00 & \centering о.е. & \centering \arraybackslash 0,01 \\
\hline
\centering Величина напряжения ОП для запрета авторегулирования (U2блк)  & \centering U2блк & \centering 0,00...2,00 & \centering о.е. & \centering \arraybackslash 0,01 \\
\hline
\centering Величина максимального напряжения для запрета авторегулирования на повышение (U\verb|>>|блк)  & \centering U\verb|>>|блк & \centering 1,00...2,00 & \centering о.е. & \centering \arraybackslash 0,01 \\
\hline
\centering Величина минимального напряжения для запрета авторегулирования (U\verb|<<|блк)  & \centering U\verb|<<|блк & \centering 0,00...1,00 & \centering о.е. & \centering \arraybackslash 0,01 \\
\hline
\centering Время задержки срабатывания на выдачу первой команды в автоматическом режиме (T1)  & \centering T1 & \centering 1000...200000 & \centering мс & \centering \arraybackslash 1 \\
\hline
\centering Время задержки срабатывания на выдачу второй и последующих команд в автоматическом режиме (T2)  & \centering T2 & \centering 100...200000 & \centering мс & \centering \arraybackslash 1 \\
\hline
\centering Время задержки выдачи ускоренной команды убавить при значительном повышении напряжения в автоматическом режиме (T3)  & \centering T3 & \centering 100...10000 & \centering мс & \centering \arraybackslash 1 \\
\hline
\centering Выдержка времени снятия команды управления после появления сигнала переключения привода (Твозвр)  & \centering Твозвр & \centering 0...10000 & \centering мс & \centering \arraybackslash 1 \\
\hline
\centering Максимальное время реакции привода РПН на команду управления (Тпуск)  & \centering Тпуск & \centering 100...10000 & \centering мс & \centering \arraybackslash 1 \\
\hline
\centering Максимальное время одного переключения привода РПН (Тперекл)  & \centering Тперекл & \centering 100...100000 & \centering мс & \centering \arraybackslash 1 \\
\hline
\centering Режим команды на отключение питания ПМ при самопроизвольном переключении (Режим\_откл\_ПМ)  & \centering Режим\_откл\_ПМ & \centering 1 секунда \\ Непрерывно (до сброса) & \centering -- & \centering \arraybackslash -- \\
\hline
\centering Ресурс привода РПН по количеству операций переключения (Ресурс\_переключений)  & \centering Ресурс\_переключений & \centering 1...1000000 & \centering -- & \centering \arraybackslash 1 \\
\hline
\centering Начальное значение счетчика переключений привода РПН (Нач\_знач\_счетчика)  & \centering Нач\_знач\_счетчика & \centering 1...1000000 & \centering -- & \centering \arraybackslash 1 \\
\hline
\centering Метод параллельного регулирования (Метод\_парал\_рег)  & \centering Метод\_парал\_рег & \centering Не используется \\ Ведущий \\ Ведомый & \centering -- & \centering \arraybackslash -- \\
\hline
\centering Контроль рассогласования приводов РПН трансформаторов в режиме параллельного регулирования (Контроль\_рассогл)  & \centering Контроль\_рассогл & \centering Внешний сигнал \\ Сигналы положений приводов РПН & \centering -- & \centering \arraybackslash -- \\
\hline
\centering Источник данных о положении РПН параллельного трансформатора (Датчик\_парал\_РПН)  & \centering Датчик\_парал\_РПН & \centering Токовая петля 4(0)-20мА \\ GOOSE \\ BCD-код & \centering -- & \centering \arraybackslash -- \\
\hline
\centering Максимально допустимый разбег приводов РПН при параллельном регулировании трансформаторов (Макс\_разбег\_РПН)  & \centering Макс\_разбег\_РПН & \centering 0...40 & \centering -- & \centering \arraybackslash 1 \\
\hline
\centering Время переключения РПН трансформаторов при параллельном регулировании (Тпарал\_перекл)  & \centering Тпарал\_перекл & \centering 100...100000 & \centering мс & \centering \arraybackslash 1 \\
\hline
\centering Номер крайнего нижнего положения РПН параллельного трансформатора (Мин\_полож\_пар\_РПН)  & \centering Мин\_полож\_пар\_РПН & \centering 1...40 & \centering -- & \centering \arraybackslash 1 \\
\hline
\centering Номер крайнего верхнего положения РПН параллельного трансформатора (Макс\_полож\_пар\_РПН)  & \centering Макс\_полож\_пар\_РПН & \centering 1...40 & \centering -- & \centering \arraybackslash 1 \\
\hline
\centering Сигнал токовой петли, соответствующий крайнему нижнему положению РПН параллельного трансформатора (Iмин\_ТП\_парал\_РПН)  & \centering Iмин\_ТП\_парал\_РПН & \centering 0...4000 & \centering мкА & \centering \arraybackslash 1 \\
\hline
\centering Сигнал токовой петли, соответствующий крайнему верхнему положению РПН параллельного трансформатора (Iмакс\_ТП\_парал\_РПН)  & \centering Iмакс\_ТП\_парал\_РПН & \centering 5000...20000 & \centering мкА & \centering \arraybackslash 1 \\
\hline
\centering Выдержка времени срабатывания сигнализации (Тсигн)  & \centering Тсигн & \centering 1...100000 & \centering мс & \centering \arraybackslash 1 \\
\hline
\centering Контроль положения блока испытательного №1 (Контроль\_БИ\_1)  & \centering Контроль\_БИ\_1 & \centering Не предусмотрено \\ Предусмотрено & \centering -- & \centering \arraybackslash -- \\
\hline
\centering Контроль положения блока испытательного №2 (Контроль\_БИ\_2)  & \centering Контроль\_БИ\_2 & \centering Не предусмотрено \\ Предусмотрено & \centering -- & \centering \arraybackslash -- \\
\hline
\centering Контроль положения блока испытательного №3 (Контроль\_БИ\_3)  & \centering Контроль\_БИ\_3 & \centering Не предусмотрено \\ Предусмотрено & \centering -- & \centering \arraybackslash -- \\
\hline
\centering Контроль положения блока испытательного №4 (Контроль\_БИ\_4)  & \centering Контроль\_БИ\_4 & \centering Не предусмотрено \\ Предусмотрено & \centering -- & \centering \arraybackslash -- \\
\hline
\centering Контроль положения блока испытательного №5 (Контроль\_БИ\_5)  & \centering Контроль\_БИ\_5 & \centering Не предусмотрено \\ Предусмотрено & \centering -- & \centering \arraybackslash -- \\
\hline
\centering Контроль положения блока испытательного №6 (Контроль\_БИ\_6)  & \centering Контроль\_БИ\_6 & \centering Не предусмотрено \\ Предусмотрено & \centering -- & \centering \arraybackslash -- \\
\hline
\centering Контроль положения блока испытательного №7 (Контроль\_БИ\_7)  & \centering Контроль\_БИ\_7 & \centering Не предусмотрено \\ Предусмотрено & \centering -- & \centering \arraybackslash -- \\
\hline
\centering Контроль положения двери (Контроль\_двери)  & \centering Контроль\_двери & \centering Не предусмотрено \\ Предусмотрено & \centering -- & \centering \arraybackslash -- \\
\hline
\centering Контроль оперативного тока технологической сигнализации (Контроль\_опер\_тока)  & \centering Контроль\_опер\_тока & \centering Не предусмотрено \\ Предусмотрено & \centering -- & \centering \arraybackslash -- \\
\hline
%===t1
\end{longtable}
\normalsize